% =============================================================================
% ch01_introduction.tex  —  Chapter 1: Introduction
% K-Bound long-form monograph.
% NO \documentclass or \usepackage. Begins with \chapter.
% =============================================================================

\chapter{Introduction}\label{ch:intro}

\section{The Deployment Problem}\label{sec:intro-deploy}

A model trained on one distribution is deployed into a world that has moved.
The training images came from one hospital's scanner; the deployed scanner is
newer, slightly sharper, with different noise characteristics.
The training network saw balanced classes; the test stream is skewed toward
one pathology.
The training corpus covered one domain; the production queries drift slowly
toward a neighbouring domain as user behaviour changes.
This pattern --- a mismatch between the distribution under which a model was
optimized and the distribution under which it operates --- is the central
operational reality of deployed machine learning.

Test-time adaptation (TTA) is the field's answer.
Rather than retraining from scratch whenever the deployment distribution
shifts, TTA updates the model on the unlabeled target data that arrives at
inference time.
The appeal is immediate: labels are expensive or unavailable in deployment,
yet the raw inputs carry information about how the distribution has moved.
Entropy minimization over batch-normalization parameters
\cite{wang2021tent}, self-supervised auxiliary objectives that continue
training at test time \cite{sun2020ttt}, source-free adaptation by
information maximization, and importance-weighted empirical risk
minimization \cite{shimodaira2000improving} are all instances of the same
basic idea: use the unlabeled stream to close the performance gap.

\section{The Double-Edged Nature of Test-Time Adaptation}\label{sec:intro-doubleedge}

Adaptation is, however, double-edged in a precise sense.
The same unlabeled objective that recovers accuracy under one kind of
distribution shift can degrade it under another --- and, critically, without
labels the deployed system cannot in general determine which situation it is in.

Consider entropy minimization.
Under mild covariate shift --- the scanner statistics change, but the
decision boundary between pathologies does not --- minimizing the entropy of
predictions sharpens the boundary toward where the new data lives.
The adaptation is helpful.
But now suppose the distribution shift is a label shift: the prevalence of
one class has collapsed to near zero in the target stream.
A model that minimizes entropy on this stream will be rewarded for making
all-or-nearly-all predictions of the dominant class; it converges to a
degenerate solution.
The adaptation is harmful, and in the worst case the model collapses
entirely.
The recent literature on TTA failure modes \cite{zhao2023pitfalls,
press2023rdumb, niu2023sar} documents exactly these collapse phenomena:
Tent's continual entropy minimization catastrophically fails on adversarial
or highly non-stationary streams; EATA's Fisher penalty slows but does not
eliminate accumulation of damage; SAR's sharpness-aware resets abort
incipient collapses but cannot prevent them in every regime.

The field's response has been to propose sturdier mechanisms.
If entropy minimization is fragile, add sample selection to filter
high-entropy inputs \cite{niu2022eata}.
If selection is not enough, add sharpness-aware updates that are less
likely to collapse \cite{niu2023sar}.
If batch normalization updates are the failure mode, use a fully
parameter-free approach \cite{boudiaf2022lame}.
If a single method is fragile, reset after detecting instability
\cite{niu2023sar} or continually re-anchor to the source
\cite{wang2022cotta}.

Each of these improvements is genuine within its scope.
But they all share a common architecture: they commit to adapting, and they
try to make that commitment safer.
None of them poses the prior question: \emph{given the available evidence,
should the system adapt at all?}

\section{The Central Question}\label{sec:intro-question}

This monograph studies that prior question.

\begin{quote}
\itshape
Can observable test-time evidence separate helpful, harmful, and
unknowable regimes of label-free adaptation, yielding a computable rule
for when a system should adapt, freeze, or abstain?
\end{quote}

The question is not rhetorical.
It has a substantive, partly negative answer that we make precise.
In some regimes --- covariate shift, certain kinds of label shift, situations
where the frozen and adapted models disagree on inputs whose label
orientation is detectable from label-free signals --- the answer is yes.
In others, the same observable evidence is compatible with both helpful and
harmful adaptation, and no label-free rule can certify the correct choice
in both worlds.
The boundary between these cases is what we call the \emph{knowability
boundary} of label-free adaptation, and the central mathematical object is
a trichotomy: at any given evidence value, adaptation is either
\emph{knowably helpful}, \emph{knowably harmful}, or \emph{unknowable}.

Why does this matter for deployment?
Consider a safety-critical setting: an autonomous vehicle's perception stack
or a medical imaging pipeline.
In such settings, the cost of silent degradation is high and the cost of
abstaining --- falling back to the frozen model and flagging for human review
--- is manageable.
A system that can identify when it is in an unknowable regime and refuse to
adapt is strictly safer than one that always adapts or always freezes,
even if its average-case performance is identical.
The value of the knowability boundary is not primarily about improving mean
accuracy; it is about converting the adapt/freeze decision from an
irreducible coin flip into a certified, controlled action.

The same logic applies to model monitoring.
A production team that can determine, from unlabeled deployment data alone,
whether a proposed update will help or hurt does not need ground-truth labels
to make a safe deployment decision.
The knowability boundary tells them exactly when that label-free determination
is possible and when it is not.

\section{The Trichotomy and Its Formalization}\label{sec:intro-trichotomy}

We formalize the trichotomy as follows.
Fix a loss $\ell$, a frozen model $f_0$, and an adapted/gated
candidate~$f_a$.
The \emph{adaptation benefit} is
\[
\Delta \;=\; R_T(f_0) - R_T(f_a),
\]
where $R_T(f) = \E_{P_T}[\ell(f(X), Y)]$ is the target risk.
A positive $\Delta$ means adapting helps; a negative $\Delta$ means it
hurts.
\emph{Observable evidence} $Z = \phi(X_{1:m}, f_0, f_a)$ is any statistic
computable without target labels: prediction entropy, softmax confidence,
score-distribution drift, detector disagreement, predicted-class balance,
update norm, or density-ratio diagnostics.
The conditional benefit $\Delta(z) = \E[\ell(f_0(X),Y) - \ell(f_a(X),Y)
\mid Z = z]$ is the benefit conditioned on evidence $Z = z$.

The regime at evidence value $z$ is:
\begin{description}
\item[Knowably helpful] if $Z$ certifies $\Delta(z) > 0$; the correct
    action is \textsc{adapt}.
\item[Knowably harmful] if $Z$ certifies $\Delta(z) < 0$; the correct
    action is \textsc{freeze}.
\item[Unknowable] if the same $z$ is compatible with both $\Delta(z) > 0$
    and $\Delta(z) < 0$; the worst-case-optimal action is
    \textsc{abstain}.
\end{description}

The key quantity is the \emph{sign} of $\Delta(z)$, not its magnitude.
Estimating the sign requires strictly less information than estimating the
risk $R_T(f)$ itself, so the knowability boundary lies strictly outside the
identifiability boundary for risk estimation --- but not everywhere: there are
situations where even the sign is unrecoverable from label-free evidence. The
rest of this chapter pins down exactly where that boundary falls, and the next
section states the three results that do it.

\section{Contributions}\label{sec:intro-contributions}

The work makes three contributions. They answer, in turn, when the
adapt/freeze question is fundamentally unanswerable from unlabeled evidence,
when it provably is answerable and how to act on it, and what happens when the
resulting rule meets real deployment streams. The manuscript opens on these
three and closes on them (Chapter~\ref{ch:discussion}); the paragraphs below
state each and the formal machinery it rests on.

\subsection*{Contribution I: The unknowable floor.}

\emph{We prove when a label-free system fundamentally cannot decide.} For any
(possibly randomized) label-free decision rule $g(Z)$, there exist two target
worlds $P_T^1$ and $P_T^2$ with identical observable-evidence laws and
opposite-sign benefits (Theorem~\ref{thm:imp}). Since the rule cannot
distinguish the worlds, its actions are identically distributed in both, and
the worst-case-optimal action is to abstain. The construction is non-vacuous:
we exhibit it explicitly as $X \sim \mathcal{N}(0,1)$ under both worlds,
$f_0(x) = \ind[x>0]$, $f_a(x) = \ind[x<0]$, and labels flipped between the
worlds. A Le~Cam two-point quantitative strengthening (Chapter~\ref{ch:theory1})
shows that the optimal committal error equals $1 - \TV$ of the evidence laws,
which is $1$ in the witness and decays to zero as the worlds become separable.
Numerically, the KGA validator confirms $100\%$ abstention on the witness and
exact tracking of the Le~Cam floor across the total-variation sweep. This is the
result that turns abstention from a heuristic fallback into a forced, principled
action.

This contribution rests on a prior formalization, developed in full in
Chapter~\ref{ch:formulation}: we cast label-free TTA as the adapt/freeze/abstain
trichotomy, define risk alignment for observable evidence, and define the three
regimes precisely. Where the TTA literature designs the adaptation
mechanism~$f_a$, we ask when the sign of $R_T(f_0) - R_T(f_a)$ is recoverable
from label-free observations at all.

\subsection*{Contribution II: The finite-sample certificate and the exact frontier.}

\emph{We prove when a label-free system provably can decide, and locate exactly
where the boundary lies.} Two results do this work. The first is operational: a
finite-sample adapt/freeze/abstain certificate with false-adapt control
(Theorem~\ref{thm:cert}). Under a label-free estimator $\widehat{\Delta}(z)$
with a calibrated radius $\varepsilon(z)$ satisfying
$\Pr[|\widehat{\Delta}(z) - \Delta(z)| \leq \varepsilon(z)] \geq 1 - \alpha$,
the rule
\[
\textsc{adapt}\ \text{if}\ \widehat{\Delta} - \varepsilon > 0,\quad
\textsc{freeze}\ \text{if}\ \widehat{\Delta} + \varepsilon < 0,\quad
\textsc{abstain}\ \text{otherwise}
\]
controls the false-adapt and false-freeze rates at $\alpha$. A corollary shows
that in the unknowable two-point regime the certificate \emph{must} abstain with
probability at least $1 - 2\alpha$, joining this contribution to the floor of
Contribution~I: where the evidence cannot tell, the certificate is forced to
abstain, not merely inclined to. A sample-complexity corollary shows the
certificate commits once $n \gtrsim (\sigma^2 / \Delta^2) \log(1/\alpha)$ paired
benefit samples are available, with an empirical-Bernstein discharge for bounded
benefits; an anytime-valid e-process variant (Appendix~\ref{app:repro}) extends
it to streaming TTA with no multiplicity correction.

The second result pins the boundary itself. Restricting to the \emph{observable
disagreement region} $D = \{x : f_0(x) \neq f_a(x)\}$ --- observable without
labels, since $f_0$ and $f_a$ are known maps of $X$ --- we prove that
$\sign\Delta$ equals an ordinal accuracy comparison on $D$: adapting helps if
and only if the adapted model exceeds chance accuracy on $D$
(Theorems~\ref{thm:disagree}--\ref{thm:disagree-reg}, for binary, multiclass,
and regression losses). This is the benefit-sign frontier: the sign is
identifiable from observables exactly when an observable margin clears the
calibration-drift budget on $D$. A companion positive regime
(Theorem~\ref{thm:pos}) shows that under covariate shift ($P_S(Y \mid X) =
P_T(Y \mid X)$) importance weighting yields a consistent target-risk estimator,
so $\sign\Delta$ is identifiable outright. Ablations confirm that detector
disagreement is the single most important evidence feature, directly supporting
the theory. Conjecture~\ref{conj:gen} states the one remaining open piece: the
label-free \emph{bracketing} of the ordinal comparison in the multiclass and
regression cases.

\subsection*{Contribution III: \KGA, a deployable safety layer.}

\emph{We make the theory a deployable rule and report its honest empirical
story.} We instantiate the framework as \KGA\ (Knowability-Guided Adaptation), a
decision wrapper placed around any TTA candidate without modifying it: it
extracts label-free evidence $Z$, estimates the benefit $\widehat{\Delta}$ with a
gradient-boosted regressor trained leave-one-task-out, forms a split-conformal
radius $\varepsilon$, and returns the three-way decision with false-adapt at most
$\alpha$. \KGA\ is released as a pure-numpy package (\texttt{kga/}) importable as
\texttt{from kga import KGA} and served at a \texttt{POST /decide} endpoint; all
code, configs, and result artifacts are public at
\url{https://github.com/Pratikn03/AutoML_Flagship_V8} \cite{kbound2026repo}.
Chapters~\ref{ch:kga} and~\ref{ch:exp1}--\ref{ch:exp2} give the algorithm and the
complete evaluation.

The empirical finding is deliberately modest and we state it plainly: \KGA\ is a
\emph{safety layer, not an accuracy booster}. It is the only policy that stays
near-oracle in both helpful and harmful regimes, which is exactly what the
trichotomy predicts --- but it earns a clean beats-both over both trivial
policies only where catastrophic, detectable harm is common. On the
pre-registered CIFAR-10-C stress grid it beats always-adapt and always-freeze for
Tent and EATA at $0\%$ false-adapt and ties collapse-resistant SAR; on
natural-shift audits it does no harm rather than winning outright. The nine
experiments below substantiate this, regime by regime.

\begin{enumerate}
\item \emph{123-task anomaly-routing benchmark} (ADBench tabular,
    image-OOD, text, cyber, fraud): \KGA\ abstains where the true benefit
    is near zero (mean $|\Delta| = 0.021$ on abstained vs.\ $0.132$ on
    acted tasks) and correctly identifies that the suite is helpful-dominated
    (adapt precision $0.90$, coverage $17\%$).
\item \emph{Harmful-regime refusal}: with a fusion candidate that hurts on
    80\% of tasks, \KGA\ refuses it almost everywhere, cutting regret-to-oracle
    by $\approx 11{\times}$ (from $0.0214$ to $0.0019$) while matching the
    safe baseline AUROC ($0.748$).
\item \emph{Controlled mixed-regime study} (369 instances, three conditions):
    \KGA\ achieves adapt precision $0.935$ and beats always-freeze by
    $0.10$ AUROC while tying always-adapt in a mild-harm domain.
\item \emph{Clean non-identifiability witness}: the exact Theorem~\ref{thm:imp}
    construction; \KGA\ abstains on $100\%$ of instances while a forced-commit
    rule pays regret $0.51$.
\item \emph{Statistical rigor over 8 seeds}: multi-seed mean AUROC
    $0.686 \pm 0.003$ with paired $t$-tests confirming \KGA\ is significantly
    above always-freeze ($p = 7.3 \times 10^{-11}$) and significantly below
    always-adapt ($p = 5.1 \times 10^{-5}$) in the mild-harm domain.
\item \emph{Regression under covariate shift}: synthetic regression with
    importance weighting hurting on 79\% of instances; \KGA\ refuses it
    entirely and matches the oracle MSE ($0.251$) vs.\ always-adapt ($0.338$).
\item \emph{Online non-stationary CIFAR-10 TTA}: across a helpful-dominated
    and a harm-dominated continual stream, \KGA\ achieves the highest mean
    accuracy across regimes ($0.490$ vs.\ always-freeze $0.456$ vs.\
    always-adapt $0.444$), with always-adapt collapsing to $0.339$ in the
    harm regime.
\item \emph{Per-corruption CIFAR-10-C helpful-dominated control} (65 cells,
    13 corruptions $\times$ 5 severities): adaptation is overwhelmingly
    helpful (false-adapt rate $0.015$); \KGA\ matches always-adapt at equal
    safety, confirming that the certificate does not impose unnecessary cost
    in benign regimes.
\item \emph{Pre-registered CIFAR-10-C stress grid} (432 conditions/method,
    harmful base rate $0.16$--$0.34$): \KGA\ beats both trivial policies for
    Tent and EATA (regret $0.0016$ vs.\ $0.009$ / $0.12$) at $0\%$
    false-adapt, and ties collapse-resistant SAR.
\end{enumerate}

\section{Honest Scope}\label{sec:intro-scope}

A safety claim is only as good as its stated limits, so we draw them in both
theory and experiment before any reader has to.

On the theory side, the impossibility result (Theorem~\ref{thm:imp}) is
close to known non-identifiability of target risk without labels; its
purpose here is to \emph{justify the abstain action} and to make the
boundary quantitative via the Le~Cam two-point lower bound.
The certificate (Theorem~\ref{thm:cert}) is proved conditional on a valid
label-free interval estimator; the empirical-Bernstein discharge requires
exchangeability of the paired benefit sample.
The sign-of-difference characterization is complete for binary, multiclass,
and regression losses; the label-free bracketing of the ordinal comparison
in the multiclass and regression cases remains an open conjecture.

On the experimental side, the stress-grid knockout (beating both trivial
policies at $0\%$ false-adapt) holds for Tent and EATA on the CIFAR-10-C
pre-registered grid; \KGA\ ties collapse-resistant SAR rather than beating
it.
In per-corruption CIFAR-10-C (helpful-dominated by construction) and within
a single-stream online regime, \KGA\ ties the better fixed policy rather
than strictly beating it.
The decisive advantage is \emph{robustness}: it is the only policy that is
near-oracle in both the helpful and harm regimes, which is exactly what
the trichotomy predicts.
Larger-scale corroboration on CIFAR-100-C, ImageNet-C, and natural-shift
benchmarks such as WILDS \cite{koh2021wilds} is future work.

The practical implication is direct.
\KGA's value scales with how often catastrophic, detectable harm occurs in
the deployment stream --- precisely the quantity the theory characterizes.
In genuinely benign domains where adaptation almost always helps, the
certificate is safety insurance at controlled cost.
In collapse-prone regimes --- small or class-imbalanced batches,
single-class label-shift streams, aggressive update schedules, long
non-stationary runs --- the certificate delivers a real win.

\section{Chapter Roadmap}\label{sec:intro-roadmap}

Chapter~\ref{ch:related} reviews the five adjacent lines of work and
positions K-Bound precisely against each.
Chapter~\ref{ch:formulation} develops the formal setup: distribution
families, the adaptation-benefit definition, observable evidence, and the
three regime definitions.
Chapters~\ref{ch:theory1} and~\ref{ch:theory2} present the five core
theoretical results.
Chapter~\ref{ch:kga} describes the \KGA\ algorithm, its assumptions, and
its cost structure.
Chapters~\ref{ch:exp1} and~\ref{ch:exp2} present the experimental
evaluation across anomaly-routing and online TTA benchmarks.
Chapter~\ref{ch:elara} presents the anonymized reliability-gated multimodal
instantiation as a fully code-validated worked example of the general
framework.
Chapter~\ref{ch:discussion} synthesizes the deployment implications,
states the limitations, specifies what would falsify the framework, and
returns to the three contributions of this chapter against the verified
evidence chain.
Appendix~\ref{app:repro} gives the full reproducibility specification.
Appendix~\ref{app:tables} contains extended tables.
Appendix~\ref{app:notation} is a notation reference.

The next chapter places K-Bound among its neighbours, building from the five
adjacent lines of work to the precise gap none of them closes.
